% !TEX root = ../main.tex

\chapter{讨论与展望}
\label{ch:discussion}

本章将对本文提出的 \textbf{SpaLineage-OT} 框架进行深入的方法学总结与讨论，剖析其在计算效率和几何表示上的物理局限性，并为未来在更大尺度时空数据上的扩展指明技术方向。

\section{研究总结}
在多细胞生物的发育生物学研究中，重建高精度的细胞发育谱系和空间迁移轨迹对于揭示基因时序调控、器官发生的物理动力学机制至关重要。本文针对传统空间最优传输模型在复杂几何拓扑下极易产生“欧氏空间瞬间移动”伪影，以及缺乏局部发育方向性约束等痛点，设计并实现了物理与几何双重约束的时空单细胞谱系追踪模型 SpaLineage-OT。

我们通过以下三个核心步骤，建立了连续时间维度的时空细胞分化流场：
\begin{enumerate}
    \item 结合 Delaunay 剖分网格与核密度估计（KDE）构建形态阻抗，在非均匀黎曼图网络上求解测地线距离，为融合格罗莫夫-瓦瑟斯坦（FGW）对齐注入了真实的组织流形结构边界约束。
    \item 将单细胞 RNA 速度向量的余弦相似度作为非对称方向性惩罚，使最优传输耦合矩阵具备了指示生物时间箭头的指向性。
    \item 基于最优传输耦合作为离散对齐桥梁，采用生成式薛定谔桥流匹配（SBFM）框架结合 DriftMLP 神经网络拟合连续时间轨迹，并在模型推断时辅以均值漂移势能负梯度力与 RK4 积分器，有效规避了细胞边缘的发散问题。
\end{enumerate}
在真实小鼠器官发育 Stereo-seq 数据集（MOSTA）上的留出验证实验证实，本模型能显著降低重建分布的能量距离（比基线 Moscot 降低了 15.5\%），并在屏障穿越率 BCR 上达到了 91.40\% 的超高水准，验证了其数学合理度与生物物理真实性。

\section{方法学讨论与瓶颈分析}
尽管 SpaLineage-OT 在 MOSTA 数据集上表现出了优异的拟合能力，但从计算生物学和大规模工业应用的视角审视，模型仍存在以下值得深入探讨的局限性：

\subsection{空间几何距离计算的复杂度瓶颈}
在 Stage 2 中，我们首先构建了基于 Delaunay 三角剖分的空间邻接图，并使用 Dijkstra 最短路径算法求解了全源最短路径距离矩阵 $D_{\text{geo}} \in \mathbb{R}^{N \times N}$。
对于拥有 $N$ 个细胞的单张切片，Dijkstra 算法的计算复杂度在配合斐波那契堆优化时为 $\mathcal{O}(N^2 \log N)$。当数据集规模处于数万细胞级别（如本研究所采用的 15,000 细胞左右的 MOSTA 切片）时，该计算耗时约为 2-3 分钟，完全在可接受范围内。然而，随着空间转录组测序分辨率的进一步提升（如百万细胞级别的新一代 Stereo-seq 切片），全源测地距离的计算和 $N \times N$ 稠密距离矩阵的内存占用将成为严重的计算瓶颈。

\subsection{静态网格拓扑对动态发育的表征局限}
本模型在 Stage 2 对每个时间点（如 E9.5 和 E11.5）独立地构建了静态的 Delaunay 空间三角剖分网络，并假设各切片内部的空间结构在当前发育周期内是静态且不发生显著宏观变形的。
然而，在真实的生物体器官发生过程中，组织不仅有细胞的游走和分化，组织本身也会发生显著的形态学伸展、折叠以及三维弯曲变形。仅依靠两个离散时间点处独立的静态网格，在应对剧烈的非线性组织形变时，可能会丢失部分动态的几何拓扑关系。

\section{未来研究展望}
为了克服上述局限，我们在未来的工作中将从以下几个技术路线对 SpaLineage-OT 进行迭代升级：

\begin{enumerate}
    \item \textbf{局部图稀疏化与 GPU/CUDA 加速}：
    针对大规模（$>10^5$ 细胞）的测地距离矩阵计算，我们将探索图稀疏化技术（如局部图划分、局部骨架提取）以降低图顶点规模。同时，设计基于 CUDA 架构的并行最短路径求解器（如并行 Bellman-Ford 或 Delta-Stepping 算法），将空间几何计算移植到 GPU 上，使计算时间缩短数倍，以支持百万级超大规模组织的时空分析。
    
    \item \textbf{动态图神经网络 (DGNN) 与流变学建模}：
    未来可引入动态图神经网络（Dynamic Graph Neural Networks）或者物理力学中的流变学（Rheological Modeling）方程，将组织骨架的形变张量也作为动力学变量融入 SBFM 模型中。这将不仅能追踪细胞个体的相对位移，还能拟合出整个胚胎组织宏观的“折叠”和“形变”场，从而在力学层面更真实地还原生命发生的规律。
    
    \item \textbf{多模态时空数据的融合建模}：
    随着单细胞空间多组学技术的迅猛发展，未来的时空数据集将不仅包含基因表达，还包含染色质可及性（Spatial ATAC-seq）、空间蛋白质组（Spatial Proteomics）以及高分辨率的细胞形态病理图像。我们将探索将高维的多组学距离度量融合进最优传输代价中，以更精确地约束细胞谱系的转录-表观-空间多维度特征演化。
\end{enumerate}
综上所述，SpaLineage-OT 将黎曼几何与连续时间薛定谔桥流匹配相结合的框架，为解析多细胞发育的动态演化提供了坚实的计算理论工具，并在时空转录组轨迹推断领域展现出了广阔的学术和实用价值。
